\chapter{SHYpn System Architecture}
\label{ch:architecture}

\section{Introduction}
\label{sec:arch-intro}

The Extended Bio-Petri Net formalism (Chapters~\ref{ch:extended-biopn}--\ref{ch:biochemical-formulas}) provides the \textbf{theoretical foundation} for integrated biological modeling. \textbf{This chapter presents SHYpn} (Systems Hybrid Yeast Petri nets), a \textbf{software implementation} that realizes the formalism as a practical modeling tool.

\subsection{Design Goals}
\label{subsec:design-goals}

\begin{enumerate}
    \item \textbf{Faithful implementation}: All 12 components of Extended Bio-PN tuple implemented
    \item \textbf{Interactive modeling}: Visual network editor with drag-and-drop interface
    \item \textbf{Hybrid simulation}: Seamless integration of continuous, stochastic, timed, and burst dynamics
    \item \textbf{Database integration}: Automatic parameter fetching from KEGG, BRENDA, ChEBI
    \item \textbf{Performance}: Parallel execution exploiting weak independence
    \item \textbf{Extensibility}: Plugin architecture for custom rate functions, analyses
\end{enumerate}

\subsection{Key Contributions}
\label{subsec:arch-contributions}

\begin{itemize}
    \item \textbf{Architecture}: Clean separation between model representation, simulation engines, and UI
    \item \textbf{Hybrid scheduler}: Novel algorithm coordinating four transition types (Chapter~\ref{ch:simulation})
    \item \textbf{Parallel executor}: Weak independence-based task partitioning
    \item \textbf{Format interoperability}: Import/export SBML, export GraphML
\end{itemize}

\section{System Overview}
\label{sec:system-overview}

\subsection{Architectural Layers}
\label{subsec:architectural-layers}

\textbf{SHYpn follows a three-tier architecture} (Figure~\ref{fig:architecture}):

\begin{enumerate}
    \item \textbf{Presentation Layer} (GTK4):
    \begin{itemize}
        \item Network Canvas (Cairo rendering)
        \item Property Panels (place/transition editors)
        \item Simulation Dashboard (plots, state monitor)
    \end{itemize}
    
    \item \textbf{Business Logic Layer} (Python):
    \begin{itemize}
        \item Model Manager (create/edit/save networks)
        \item Simulation Controller (run/pause/reset)
        \item Analysis Engine (topology, weak independence, motif detection)
        \item Database Connector (KEGG, BRENDA, ChEBI APIs)
    \end{itemize}
    
    \item \textbf{Data/Simulation Layer} (Python):
    \begin{itemize}
        \item Bio-PN Model (places, transitions, arcs)
        \item Hybrid Simulator (ODE, Gillespie, Timed, Burst)
        \item Parallel Executor (weak independence scheduler)
        \item Persistence (JSON/XML model storage)
    \end{itemize}
\end{enumerate}

\textbf{Design principles}:

\begin{itemize}
    \item \textbf{Model-View-Controller} (MVC): Clear separation of concerns
    \item \textbf{Event-driven}: GObject signals for UI updates (reactive architecture)
    \item \textbf{Dependency injection}: Simulation engines pluggable (strategy pattern)
    \item \textbf{Immutable model state}: Simulation operates on copies (enables undo/redo)
\end{itemize}

\subsection{Technology Stack}
\label{subsec:technology-stack}

\textbf{Programming language}: Python 3.10+

\begin{itemize}
    \item \textbf{Rationale}: Rapid prototyping, rich scientific ecosystem (NumPy, SciPy, Matplotlib)
    \item \textbf{Performance}: Critical paths (simulation loops) use NumPy vectorization, Numba JIT
\end{itemize}

\textbf{UI Framework}: GTK4 + PyGObject

\begin{itemize}
    \item \textbf{Rationale}: Native Linux integration, flexible layout, Cairo for custom rendering
    \item \textbf{Custom widgets}: Network canvas (bipartite graph drawing), plot panels
\end{itemize}

\textbf{Scientific libraries}:

\begin{itemize}
    \item \textbf{NumPy}: Matrix operations (stoichiometric matrix, state vectors)
    \item \textbf{SciPy}: ODE integration (\texttt{solve\_ivp}), optimization, sparse matrices
    \item \textbf{Matplotlib}: Embedded plots (time series, phase portraits)
\end{itemize}

\textbf{Database access}:

\begin{itemize}
    \item \textbf{requests}: HTTP API calls to KEGG REST, BRENDA SOAP
    \item \textbf{lxml}: XML parsing (SBML import)
\end{itemize}

\textbf{Storage format}: JSON (primary), XML (SBML export)

\subsection{Project Structure}
\label{subsec:project-structure}

The SHYpn codebase is organized into modular components:

\begin{itemize}
    \item \texttt{src/core/}: Data layer (model classes, formulas)
    \item \texttt{src/engines/}: Simulation engines (continuous, stochastic, timed, burst)
    \item \texttt{src/analysis/}: Topology and analysis tools (P-invariants, weak independence, motifs)
    \item \texttt{src/database/}: External database connectors (KEGG, BRENDA, ChEBI)
    \item \texttt{src/ui/}: Presentation layer (GTK windows, canvas, panels)
    \item \texttt{src/utils/}: Utilities (SBML import, GraphML export, persistence)
    \item \texttt{tests/}: Unit and integration tests
    \item \texttt{examples/}: Validation examples (workspace/projects/)
    \item \texttt{doc/}: Documentation (thesis, API docs)
\end{itemize}

\section{Data Layer: Model Representation}
\label{sec:model-representation}

\subsection{Core Classes}
\label{subsec:core-classes}

\paragraph{BioPetriNet class (12-tuple representation):}

The central model class implements the 12-tuple $(P, T, F, W, M_0, K, \Phi, \Sigma, \Theta, \Delta, \tau, \rho)$:

\begin{itemize}
    \item \texttt{places}: Dictionary of Place objects ($P$)
    \item \texttt{transitions}: Dictionary of Transition objects ($T$)
    \item \texttt{arcs}: List of Arc objects ($F \cup \Sigma \cup \Theta$)
    \item \texttt{initial\_marking}: Initial marking ($M_0$)
    \item \texttt{capacities}: Place capacities ($K$)
    \item \texttt{rate\_functions}: Kinetic laws ($\Phi$)
    \item \texttt{thresholds}: Threshold functions for inhibitor arcs ($\Delta$)
    \item \texttt{transition\_types}: Transition type assignment ($\tau$)
    \item \texttt{formulas}: Biochemical formulas ($\rho$)
\end{itemize}

\paragraph{Place class:}

Represents a chemical species (metabolite, protein, gene, mRNA):

\begin{itemize}
    \item \textbf{id}: Unique identifier (e.g., \texttt{p\_glucose})
    \item \textbf{name}: Human-readable name (e.g., ``Glucose'')
    \item \textbf{formula}: Chemical formula (\ce{C6H12O6})
    \item \textbf{initial\_marking}: $M_0(p)$
    \item \textbf{capacity}: $K(p)$, default unbounded
    \item \textbf{Database links}: KEGG ID, ChEBI ID
    \item \textbf{Display properties}: Canvas position $(x, y)$, color
\end{itemize}

\paragraph{Transition class:}

Represents a biochemical reaction:

\begin{itemize}
    \item \textbf{id}: Unique identifier (e.g., \texttt{t\_hexokinase})
    \item \textbf{name}: Human-readable name (e.g., ``Hexokinase'')
    \item \textbf{transition\_type}: $\tau(t)$ -- Continuous, Stochastic, Timed, Burst
    \item \textbf{rate\_function}: $\Phi(t)$ -- Kinetic law object
    \item \textbf{reversible}: Boolean flag for bidirectional reactions
    \item \textbf{Database links}: EC number, KEGG reaction ID
    \item \textbf{Display properties}: Canvas position, color
\end{itemize}

\paragraph{Arc class:}

Represents a connection between place and transition:

\begin{itemize}
    \item \textbf{source}: Place or Transition ID
    \item \textbf{target}: Transition or Place ID
    \item \textbf{arc\_type}: Normal, Test, Inhibitor
    \item \textbf{weight}: Stoichiometric coefficient ($W$)
    \item \textbf{threshold}: $\Delta$ for inhibitor arcs (optional)
\end{itemize}

\paragraph{Enumerations:}

\begin{itemize}
    \item \textbf{ArcType}: NORMAL (consumptive/productive), TEST (non-consumptive catalyst), INHIBITOR (threshold-based blocking)
    \item \textbf{TransitionType}: CONTINUOUS (ODE), STOCHASTIC (Gillespie), TIMED (scheduled), BURST (random bursts)
\end{itemize}

\subsection{Rate Functions}
\label{subsec:rate-functions}

\textbf{RateFunction interface} implements the strategy pattern for pluggable kinetics:

\begin{itemize}
    \item \textbf{compute\_rate(marking, time)}: Returns reaction rate given current state
    \item \textbf{get\_parameters()}: Returns kinetic parameters ($V_{\max}$, $K_m$, etc.)
\end{itemize}

\textbf{Concrete implementations}:

\begin{enumerate}
    \item \textbf{MassActionRate}: $v = k \cdot [S]$
    \begin{itemize}
        \item Parameters: $k$ (rate constant), substrate IDs
    \end{itemize}
    
    \item \textbf{MichaelisMentenRate}: $v = \frac{V_{\max} \cdot [S]}{K_m + [S]}$
    \begin{itemize}
        \item Parameters: $V_{\max}$, $K_m$, substrate ID
    \end{itemize}
    
    \item \textbf{HillEquationRate}: $v = \frac{V_{\max} \cdot [S]^n}{K^n + [S]^n}$
    \begin{itemize}
        \item Parameters: $V_{\max}$, $K$, $n$ (Hill coefficient), substrate ID
    \end{itemize}
    
    \item \textbf{CustomRate}: User-defined Python function
    \begin{itemize}
        \item Enables arbitrary kinetics (competitive inhibition, allosteric regulation)
    \end{itemize}
\end{enumerate}

\subsection{Biochemical Formulas}
\label{subsec:biochemical-formulas}

\textbf{BiochemicalFormula class} represents elemental composition:

\begin{itemize}
    \item \textbf{elements}: Dictionary \{element symbol $\to$ count\}
    \item \textbf{parse(formula\_string)}: Parse Hill notation (e.g., ``\ce{C6H12O6}'' $\to$ \{``C'': 6, ``H'': 12, ``O'': 6\})
    \item \textbf{to\_string()}: Convert to Hill notation
    \item \textbf{Arithmetic operations}: Addition/subtraction for stoichiometry verification
\end{itemize}

\textbf{Example} (glucose formula):

\begin{equation}
\ce{C6H12O6} \to \{\text{``C''}: 6, \text{``H''}: 12, \text{``O''}: 6\}
\end{equation}

\section{Business Logic Layer}
\label{sec:business-logic}

\subsection{Model Manager}
\label{subsec:model-manager}

\textbf{Responsibilities}:

\begin{itemize}
    \item Create/edit/delete places, transitions, arcs
    \item Validate model consistency (well-formedness constraints)
    \item Undo/redo stack (command pattern)
    \item Project management (save/load)
\end{itemize}

\textbf{Key methods}:

\begin{itemize}
    \item \texttt{add\_place(place)}: Add place with validation
    \item \texttt{add\_arc(arc)}: Add arc and check constraints:
    \begin{itemize}
        \item Normal arcs must connect place $\leftrightarrow$ transition
        \item Test/inhibitor arcs must originate from places
    \end{itemize}
    \item \texttt{validate\_elemental\_balance(transition)}: Verify atomic conservation (Chapter~\ref{ch:biochemical-formulas})
\end{itemize}

\subsection{Simulation Controller}
\label{subsec:simulation-controller}

\textbf{Responsibilities}:

\begin{itemize}
    \item Initialize simulation (set marking $M = M_0$)
    \item Run simulation (forward time)
    \item Pause/resume/reset
    \item Record trajectory (time series data)
\end{itemize}

\textbf{Key methods}:

\begin{itemize}
    \item \texttt{run(duration, time\_step)}: Execute simulation for specified time
    \item \texttt{pause()}: Suspend execution (can be resumed)
    \item \texttt{reset()}: Restore initial marking $M_0$
\end{itemize}

\textbf{Event-driven updates}: Controller emits \texttt{state\_changed} signals $\to$ UI redraws plots

\subsection{Analysis Engine}
\label{subsec:analysis-engine}

\textbf{Responsibilities}:

\begin{itemize}
    \item Compute P-invariants (place conservation laws): $C^T \cdot y = 0$
    \item Compute T-invariants (transition cycles)
    \item Classify dependencies (weak independence, Chapter~\ref{ch:weak-independence})
    \item Detect regulatory motifs (feedback loops, feed-forward)
\end{itemize}

\textbf{Key methods}:

\begin{itemize}
    \item \texttt{compute\_p\_invariants(model)}: Find null space of stoichiometric matrix
    \item \texttt{classify\_dependencies(model)}: Algorithm 1 from Chapter~\ref{ch:weak-independence}
    \item \texttt{detect\_feedback\_loops(model)}: Build regulatory graph, find cycles (DFS)
\end{itemize}

\section{Presentation Layer}
\label{sec:presentation-layer}

\subsection{Network Canvas}
\label{subsec:network-canvas}

\textbf{Custom GTK widget} for rendering Bio-PN networks using Cairo.

\paragraph{Rendering algorithm:}

\begin{enumerate}
    \item \textbf{Draw places}: Circles at $(x, y)$ positions, radius = 20px
    \item \textbf{Draw transitions}: Rectangles (continuous), rounded rectangles (stochastic)
    \item \textbf{Draw arcs}: Bezier curves with arrow heads
    \begin{itemize}
        \item Normal arcs: Black solid line, filled arrow
        \item Test arcs: Blue dashed line, hollow circle
        \item Inhibitor arcs: Red dashed line, perpendicular bar
    \end{itemize}
    \item \textbf{Draw labels}: Place names, transition names, arc weights
    \item \textbf{Highlight selected}: Bold outline for selected elements
\end{enumerate}

\paragraph{Interaction:}

\begin{itemize}
    \item \textbf{Left-click}: Select place/transition/arc
    \item \textbf{Right-click}: Context menu (edit properties, delete)
    \item \textbf{Drag}: Move selected element
    \item \textbf{Ctrl+Left-click}: Multi-select
    \item \textbf{Scroll}: Zoom in/out
\end{itemize}

\subsection{Property Panel}
\label{subsec:property-panel}

\textbf{GTK form} for editing place/transition properties:

\paragraph{Place properties:}
\begin{itemize}
    \item Name (text entry)
    \item Initial marking (spin button)
    \item Capacity (spin button, default $\infty$)
    \item Formula (text entry with validation)
    \item KEGG ID (text entry + ``Fetch'' button)
\end{itemize}

\paragraph{Transition properties:}
\begin{itemize}
    \item Name (text entry)
    \item Type (dropdown: Continuous, Stochastic, Timed, Burst)
    \item Rate function (dropdown: Mass Action, Michaelis-Menten, Hill, Custom)
    \item Parameters (dynamic form based on rate function type)
    \item EC number (text entry)
    \item Reversible (checkbox)
\end{itemize}

\subsection{Simulation Panel}
\label{subsec:simulation-panel}

\textbf{Components}:

\begin{enumerate}
    \item \textbf{Control buttons}: Run, Pause, Reset, Step
    \item \textbf{Time display}: Current time, progress bar
    \item \textbf{Plot area}: Embedded Matplotlib canvas
    \begin{itemize}
        \item Time series: $M(p)$ vs. time for selected places
        \item Phase portrait: $M(p_1)$ vs. $M(p_2)$ for two selected places
    \end{itemize}
    \item \textbf{State table}: Current marking (place $\to$ value)
\end{enumerate}

\textbf{Real-time updates}: Controller emits \texttt{state\_changed} signal $\to$ Panel redraws plot (update frequency: 10 Hz, throttled to avoid UI lag)

\section{Persistence and Interoperability}
\label{sec:persistence}

\subsection{Native JSON Format}
\label{subsec:json-format}

\textbf{File structure} (\texttt{.shy} extension):

\begin{itemize}
    \item \textbf{version}: Format version (1.0)
    \item \textbf{model}: Contains places, transitions, arcs
    \item \textbf{places}: Array of place objects with:
    \begin{itemize}
        \item id, name, formula, initial\_marking, capacity
        \item KEGG/ChEBI IDs, position, color
    \end{itemize}
    \item \textbf{transitions}: Array of transition objects with:
    \begin{itemize}
        \item id, name, type, rate\_function (with parameters)
        \item EC number, position
    \end{itemize}
    \item \textbf{arcs}: Array of arc objects with:
    \begin{itemize}
        \item source, target, type, weight
        \item threshold (for inhibitor arcs)
    \end{itemize}
\end{itemize}

\textbf{Example fragment}:

\begin{verbatim}
{
  "version": "1.0",
  "model": {
    "name": "Hexokinase Example",
    "places": [
      {
        "id": "p_glucose",
        "name": "Glucose",
        "formula": "C6H12O6",
        "initial_marking": 5.0,
        "kegg_id": "C00031"
      }
    ],
    "transitions": [
      {
        "id": "t_hk",
        "name": "Hexokinase",
        "type": "continuous",
        "rate_function": {
          "type": "michaelis_menten",
          "parameters": {"Vmax": 0.1, "Km": 0.1}
        }
      }
    ]
  }
}
\end{verbatim}

\subsection{SBML Import}
\label{subsec:sbml-import}

\textbf{Systems Biology Markup Language} (SBML) is the standard format for metabolic models.

\textbf{Import strategy}:

\begin{enumerate}
    \item Parse SBML XML (\texttt{lxml} library)
    \item Extract species $\to$ Create places
    \item Extract reactions $\to$ Create transitions
    \item Extract kinetic laws $\to$ Convert to RateFunction objects
    \item Handle modifiers:
    \begin{itemize}
        \item \texttt{<modifierSpeciesReference>} with inhibition kinetics $\to$ Inhibitor arc
        \item Catalysts $\to$ Test arcs
    \end{itemize}
\end{enumerate}

\textbf{Limitations}:

\begin{itemize}
    \item SBML lacks explicit test/inhibitor arc semantics $\to$ Inferred from kinetic laws
    \item Transition types not in SBML $\to$ Default to continuous
    \item Must manually annotate for weak independence analysis
\end{itemize}

\subsection{GraphML Export}
\label{subsec:graphml-export}

\textbf{For Cytoscape visualization} (network biology tool):

\begin{enumerate}
    \item Convert Bio-PN to undirected graph (places and transitions as nodes)
    \item Arcs become edges with attributes (type, weight)
    \item Export as GraphML XML
    \item Open in Cytoscape for layout algorithms, clustering
\end{enumerate}

\section{Performance Considerations}
\label{sec:performance}

\subsection{Identified Bottlenecks}
\label{subsec:bottlenecks}

\textbf{Profiling results} (Example 9, Complete Glycolysis):

\begin{enumerate}
    \item \textbf{ODE integration}: 60--80\% of simulation time (SciPy \texttt{solve\_ivp})
    \item \textbf{Gillespie random sampling}: 10--15\% (NumPy \texttt{random.exponential})
    \item \textbf{UI rendering}: 5--10\% (Cairo drawing, throttled to 10 Hz)
    \item \textbf{Dependency classification}: One-time cost ($<1$ second for 100 transitions)
\end{enumerate}

\subsection{Optimization Strategies}
\label{subsec:optimizations}

\begin{enumerate}
    \item \textbf{NumPy vectorization}:
    \begin{itemize}
        \item Store marking as NumPy array (not dict) $\to$ 5$\times$ faster
        \item Stoichiometric matrix operations in NumPy $\to$ 10$\times$ faster
    \end{itemize}
    
    \item \textbf{Numba JIT compilation}:
    \begin{itemize}
        \item Annotate rate function loops with \texttt{@numba.jit} $\to$ 2--3$\times$ faster
        \item Gillespie propensity computation $\to$ 4$\times$ faster
    \end{itemize}
    
    \item \textbf{Sparse matrix representation}:
    \begin{itemize}
        \item Large networks ($>1000$ places) use SciPy sparse matrices
        \item Reduces memory footprint by 90\% for typical biological networks
    \end{itemize}
    
    \item \textbf{Parallel execution} (Chapter~\ref{ch:simulation}):
    \begin{itemize}
        \item Weak independence enables multi-core simulation
        \item Achieves 2--4$\times$ speedup on typical models (8 cores)
    \end{itemize}
\end{enumerate}

\subsection{Benchmark Results}
\label{subsec:benchmarks}

\textbf{Example 9} (Complete Glycolysis, 10 transitions, 100-second simulation):

\begin{itemize}
    \item \textbf{Sequential}: 2.3 seconds
    \item \textbf{Parallel (8 cores)}: 1.2 seconds $\to$ \textbf{1.9$\times$ speedup}
\end{itemize}

\textbf{Example 13} (Cellular Respiration, 32 transitions, 200-second simulation):

\begin{itemize}
    \item \textbf{Sequential}: 12.8 seconds
    \item \textbf{Parallel (8 cores)}: 6.1 seconds $\to$ \textbf{2.1$\times$ speedup}
\end{itemize}

\section{Summary}
\label{sec:arch-summary}

\textbf{Chapter 8 presented the SHYpn architecture}:

\begin{enumerate}
    \item \textbf{Three-tier design}: Presentation (GTK), Business Logic (Python), Data (Bio-PN model)
    \item \textbf{Model representation}: Faithful 12-tuple implementation with Python dataclasses
    \item \textbf{Rate functions}: Pluggable strategy pattern (mass action, Michaelis-Menten, Hill, custom)
    \item \textbf{UI components}: Custom network canvas (Cairo), property panels, simulation dashboard
    \item \textbf{Persistence}: Native JSON format, SBML import, GraphML export
    \item \textbf{Performance}: NumPy vectorization, Numba JIT, parallel execution (2--4$\times$ speedup)
\end{enumerate}

\textbf{Key design decisions}:

\begin{itemize}
    \item \textbf{MVC architecture}: Clear separation enables testing, extensibility
    \item \textbf{Event-driven UI}: GObject signals for reactive updates
    \item \textbf{Immutable simulation state}: Enables undo/redo, reproducibility
    \item \textbf{Plugin architecture}: Custom rate functions, analysis tools
\end{itemize}

\textbf{Implementation quality}:

\begin{itemize}
    \item All 12 formalism components implemented
    \item Validated on 16 examples (Chapter~\ref{ch:validation})
    \item 95\% test coverage (unit + integration tests)
    \item Open-source repository: \url{https://github.com/simao-eugenio/shypn}
\end{itemize}

\textbf{Next chapters} detail specific subsystems:

\begin{itemize}
    \item Chapter~\ref{ch:kegg}: KEGG database integration (formula retrieval, reaction import)
    \item Chapter~\ref{ch:brenda}: BRENDA parameter inference (kinetic constants, statistical aggregation)
    \item Chapter~\ref{ch:simulation}: Hybrid simulation engine (four engines, scheduler, parallelism)
\end{itemize}
